\begin{figure}[ht]
\centering
\begin{tikzpicture}[x=1cm, y=1cm, font=\small]

% Euler-Maclaurin: the sum of f(n) as boxes of unit width against the
% smooth integral of f(x). f(x) = 4/(x+1) on [1,7], a decreasing convex
% curve. The rectangles (left endpoints) overshoot the integral; the
% shaded slivers between rectangle tops and the curve are the endpoint
% and Bernoulli corrections the formula accounts for exactly.

% axes
\draw[->, thick] (-0.2, 0) -- (8.6, 0) node[right] {$x$};
\draw[->, thick] (0, -0.2) -- (0, 4.4) node[above] {$f(x)$};

% x ticks (integer nodes n = 1..7)
\foreach \n in {1,2,3,4,5,6,7} {
  \draw (\n, -0.08) -- (\n, 0.08) node[below, yshift=-0.12cm] {\n};
}
% y ticks
\foreach \y in {1,2,3} {
  \draw (-0.08, \y) -- (0.08, \y) node[left, xshift=-0.1cm] {\y};
}

% Left-endpoint rectangles f(n) over [n, n+1], f(x)=4/(x+1): f(1)=2, f(2)=1.333,
% f(3)=1, f(4)=0.8, f(5)=0.667, f(6)=0.571
\foreach \n/\h in {1/2.0, 2/1.3333, 3/1.0, 4/0.8, 5/0.6667, 6/0.5714} {
  \draw[enggray, fill=engblue!12] (\n, 0) rectangle ({\n+1}, \h);
}

% Smooth curve f(x) = 4/(x+1) on [1,7]
\draw[very thick, engred] plot[domain=1:7, samples=80] (\x, {4/(\x+1)})
  node[right] {$f(x)$};

% mark the sample heights at nodes
\foreach \n/\h in {1/2.0, 2/1.3333, 3/1.0, 4/0.8, 5/0.6667, 6/0.5714} {
  \fill[engblue] (\n, \h) circle (1.2pt);
}

% annotation: the gap slivers
\draw[->, engblack] (5.9, 2.6) -- (2.55, 1.3);
\node[engblack, anchor=west, align=left] at (5.85, 2.85)
  {\footnotesize sum of boxes\\[-1pt]\footnotesize $-$ integral $=$\\[-1pt]\footnotesize
   endpoint $+$ Bernoulli\\[-1pt]\footnotesize corrections};

\end{tikzpicture}
\caption[The Euler-Maclaurin gap between a sum and its integral]{The
Euler-Maclaurin formula as the exact accounting of the gap between a
sum and an integral. The rectangles are the terms $f(n)$ read as unit-
width blocks; the smooth curve is $\int f$. For a decreasing convex $f$
the left-endpoint blocks overshoot the area under the curve, and the
total overshoot is not a vague inequality but a closed sum: the
trapezoidal endpoint average $\tfrac{1}{2}(f(a)+f(b))$ plus the
Bernoulli-number derivative corrections. Where the integral test
returns only the bound that the sum exceeds the integral, the
Euler-Maclaurin formula returns the overshoot term by term, which is
why a modest number of evaluated terms plus a closed tail beats
grinding through thousands of terms of a slowly convergent series.}
\label{fig:vol02:ch04:euler-maclaurin}
\end{figure}
